\documentclass[11pt]{article}
\usepackage[T1]{fontenc}
\usepackage[utf8]{inputenc}
\usepackage{lmodern,microtype}
\usepackage[margin=1in]{geometry}
\usepackage{amsmath,amssymb,amsthm,mathtools}
\usepackage{booktabs,array,longtable}
\usepackage{xcolor,listings,fancyhdr}
\usepackage[colorlinks=true,linkcolor=blue!45!black,urlcolor=blue!45!black,citecolor=blue!45!black]{hyperref}
\hypersetup{pdftitle={A smaller induction selector for iPRA},
  pdfauthor={Research note prepared for Vladimir Reshetnikov},
  pdfsubject={Unrestricted mirror-selected six-rule calculus and bi-interpretation with iPRA}}
\usepackage{enumitem}
\setlist{itemsep=3pt,topsep=5pt}
\setlength{\emergencystretch}{2em}
\setlength{\headheight}{14pt}
\pagestyle{fancy}\fancyhf{}
\fancyhead[L]{\small A smaller induction selector for iPRA}
\fancyhead[R]{\small Mirror core}
\fancyfoot[C]{\thepage}
\newtheorem{theorem}{Theorem}[section]
\newtheorem{lemma}[theorem]{Lemma}
\newtheorem{proposition}[theorem]{Proposition}
\newtheorem{corollary}[theorem]{Corollary}
\theoremstyle{definition}\newtheorem{definition}[theorem]{Definition}
\newcommand{\Mcal}{\mathcal M}
\newcommand{\Ncal}{\mathcal N}
\newcommand{\Scal}{\mathcal S}
\newcommand{\FV}{\operatorname{FV}}
\newcommand{\Mir}{\operatorname{Mir}}
\newcommand{\ipra}{\mathsf{iPRA}}
\newcommand{\ceq}{\mathrel{\equiv_{\rm syn}}}
\lstset{basicstyle=\ttfamily\footnotesize,columns=fullflexible,
  keepspaces=true,breaklines=true,showstringspaces=false,
  frame=single,rulecolor=\color{black!15},xleftmargin=4pt,xrightmargin=4pt,
  keywordstyle=\color{blue!45!black},commentstyle=\color{black!55}}
\title{A smaller induction selector for iPRA\\
  \large An unrestricted six-rule presentation with one mirrored iterator}
\author{Research note prepared for Vladimir Reshetnikov}
\date{November 11, 2026}
\begin{document}
\maketitle
\begin{abstract}
The current ConciseIPRA core already has six rule families on raw iterator
terms and unrestricted equality/$Q$/$D$ formulas. This note makes a modest
further simplification: replace its two-iterator diagonal test for
progressive introduction by the one-iterator pattern
$a(x,y)=I(y,a(x,y);u.h)$. The language, ordinary readings, and remaining
rules are unchanged. We prove equality of the complete contextual
consequence relations, and hence transfer both interpretation composites
and equiconsistency with full first-order intuitionistic primitive
recursive arithmetic. Arbitrary numerical substitution is admissible by
direct induction, using an inherited repair for newly selected bases.
Matched standalone Python checkers differ only in the selector: the
normalized source decreases from 3,125 to 3,109 bytes, from 1,022 to 1,012
significant tokens, and from 1,091 to 1,082 AST nodes. The gain is a
specialization and combination of existing ideas, not a new rule-count
reduction, global checker-size record, or absolute minimality theorem.
\end{abstract}
\tableofcontents
\newpage

\section{Result, baseline, and scope}

The comparison is pinned to ProveIt commit
\texttt{e1d9e3b12edd642ea2f5b572bb43febceaa50d93}, in
\texttt{Logic/PrimitiveRecursiveArithmetic/ConciseIPRA}. The current
README indexes 67 research packages and a Lean development. In particular,
\texttt{NativeSix.lean} defines the six-family calculus \texttt{SPrf},
proves \texttt{sprf\_iff\_nprf}, and proves substitution admissibility
\cite{readme,six}. That is the mathematical baseline used here.

\begin{theorem}[Main result]\label{thm:main}
Let $\Mcal$ be the calculus in Section~\ref{sec:spec}, and let $\Ncal$
be the current eight-constructor native kernel \texttt{NPrf}.
For every finite context $\Gamma$ and every native formula $A$,
\[
 \Gamma\vdash_{\Mcal}A
 \quad\Longleftrightarrow\quad
 \Gamma\vdash_{\Ncal}A.
\]
Both transformations are effective on finite derivations and leave the
sequent syntax unchanged. Consequently $\Mcal$ presents a theory
bi-interpretable and equiconsistent with full first-order $\ipra$.
Both numerical interpretation domains are whole; equality is actual
equality; and the two numerical composite isomorphisms have graph $x=y$.
\end{theorem}

The precise simplification is in the introduction selector:
\[
\begin{array}{ll}
\text{Current diagonal:}&I(y,y;u.h)=I(y,y;v.k),\\[2pt]
\text{Proposed mirror:}&a(x,y)=I(y,a(x,y);u.h).
\end{array}
\]
The latter recognizes one iterator skeleton and a repeated term. It also
allows a one-iterator lift of an arbitrary equation when recovering
induction. It does not restrict the grammar of public or intermediate
formulas. Neither a normalizer nor a primitive term-template parameter
is added to the checker.

\paragraph{Attribution.}
The mirror shape is already present in \texttt{Research/mirror\_guard},
where it is part of a restricted formula grammar and a different
transport interface \cite{mirror}. Two-case Eval is already present in
the current Lean development \cite{six}. The observation that a guard
need only be preserved under substitution, with newly true cases repaired
by a provable base, is explicitly established in
\texttt{Research/reflexive\_base/proof.tex}, including a raw-syntax
variant \cite{reflexive}. Those ingredients are inherited. The present
contribution is their particular combination on the full raw grammar,
the exact comparison with \texttt{SPrf}, and the supplied executable
evidence. No priority beyond the inspected snapshot is claimed.

\section{Complete specification of the proposed calculus}\label{sec:spec}

\subsection{Numerical terms and formulas}
There is one numerical sort. The raw term grammar is
\[
 t ::= x\mid 0\mid S(t)\mid I(n,a;u.h).
\]
Only the state variable $u$ is bound in $h$. The counter $n$ and
initializer $a$ are outside that binder. An update can have fixed
numerical parameters. Terms are finite syntax trees, not equivalence
classes. Substitution is ordinary capture-free substitution, with no
normalization step.

The formula grammar is unrestricted:
\[
 A ::= (s=t)\mid Q_y A\mid D_{x,y}A.
\]
The variables bound by $D$ are distinct. Write $B(t,r)$ for simultaneous
capture-free replacement of its distinguished free coordinates $x,y$ by
$t,r$. The designated contradiction is $0=S0$. There are no primitive
ordinary connectives, falsity symbol, predicate symbols other than
equality, or additional numerical operations.

For orientation, the ordinary readings are
\[
 Q_yB(y) := B(0)\to\forall y\,B(y),
 \qquad
 D_{x,y}B(x,y) :=
 \exists x\bigl(B(x,0)\to\forall y\,B(x,y)\bigr).
\]
These are interpretation clauses, not extra axioms of the native
calculus. Contexts are finite lists of reusable formulas. Their order
and multiplicities have no logical effect, although a checker may use
indices to identify a hypothesis.

\subsection{The selector}
Define $\Mir(B)$ to hold exactly when the \emph{uninstantiated} matrix
$B(x,y)$ is literally
\[
             a(x,y)=I(y,a(x,y);u.h(x,y,u)).             \tag{M}
\]
The counter must be the distinguished bound variable $y$. The
initializer must be the same raw term tree as the left side of the
equation. The body $h$ is arbitrary, and its state binder is fresh.
Define
\[
 \eta_B(z)=\begin{cases}z,&\Mir(B),\\0,&\neg\Mir(B).\end{cases}
\]
This is a syntactic test. It asks for neither arithmetic equality nor
provability. It is performed before substituting the existential
witness. For example $x=I(y,0;u.0)$ is not selected, even when the
subsequently supplied witness is $0$.

\subsection{All primitive rules}
The following six families are the entire calculus:
\begin{align*}
 &\frac{A\in\Gamma}{\Gamma\vdash A}
       &&\mathrm{Hyp},\\[5pt]
 &\frac{}{\Gamma\vdash t=t}
       &&\mathrm{Refl},\\[5pt]
 &\frac{\Gamma\vdash a=t\qquad
         \Gamma\vdash I(0,a;u.h)=r}
        {\Gamma\vdash r=t}
       &&\mathrm{Eval}_0,\\[5pt]
 &\frac{\Gamma\vdash I(c,a;u.h)=t\qquad
         \Gamma\vdash I(Sc,a;u.h)=r}
        {\Gamma\vdash r=h[t/u]}
       &&\mathrm{Eval}_S,\\[5pt]
 &\frac{\Gamma\vdash Q_yB(y)\qquad\Gamma\vdash B(0)}
        {\Gamma\vdash B(t)}
       &&\mathrm{QE},\\[5pt]
 &\frac{\Gamma\vdash D_{x,y}B(x,y)\qquad
         \Gamma,Q_yB(b,y)\vdash A}
        {\Gamma\vdash A}
       &&\mathrm{DE},\\[5pt]
 &\frac{\Gamma,B(t,\eta_B(z))\vdash B(t,Sz)}
        {\Gamma\vdash D_{x,y}B(x,y)}
       &&\mathrm{DI}_{M}.
\end{align*}
The two Eval rows are the two cases of one family, exactly as in the
baseline's six-rule count. There are seven inference constructors if
those cases are counted separately. This convention is explicit in both
systems; no rule has been hidden by a change of counting convention.

In DE, $b$ is fresh for $\Gamma$, $A$, and every passive parameter of
$B$. In DI, $z$ is fresh for $\Gamma$, $t$, and every passive parameter
of $B$. The displayed local assumptions are discharged. All substitutions
avoid capture. There are no other logical or nonlogical axioms, no
primitive structural rules, no successor axioms, no separate computation
axioms, and no general induction rule. All of these needed operations
are derived below.

\section{Structural operations, equality, and computation}

\subsection{The order of structural arguments}
\begin{lemma}\label{lem:struct}
Weakening, injective renaming of numerical variables, and simultaneous
ordinary-assumption cut are admissible in $\Mcal$, without using
general numerical substitution or equation induction.
\end{lemma}
\begin{proof}
Weakening is induction on derivations. An injective renaming preserves
constructor shapes and is injective on raw term trees. It therefore
preserves and reflects both the distinguished counter test and equality
of the two mirrored occurrences of $a$. It preserves every primitive
rule, including the DI mode. Renaming a derivation into a larger fresh
scope is a special case.

For cut, induct on the derivation into which the supplied proofs are
inserted. Replace a hypothesis leaf by its supplied derivation. At a
binder rule, rename the eigenvariable away from the finitely many
inserted data, lift the supplied derivations by the injective renaming
just proved, and retain the local assumption by Hyp. The matrix of DI
is unaffected by replacement of hypotheses, so its mode is unchanged.
This proves simultaneous cut and its one-assumption specialization.
\end{proof}

Exchange and contraction are immediate consequences of reusable
hypotheses and weakening. General numerical substitution is deliberately
postponed: a noninjective substitution can make a previously false
mirror test true.

\subsection{Explicit equality derivations}
All the following arguments use Refl and Eval alone, so they are the
same elementary bootstrap as in \texttt{NativeSix} \cite{six}.

\begin{lemma}[Joining equations]\label{lem:join}
From $s=t$ and $s=r$ derive $r=t$.
\end{lemma}
\begin{proof}
Put $J=I(0,s;u.u)$. Eval$_0$ on $s=r$ and $J=J$ gives $J=r$.
Eval$_0$ on $s=t$ and this $J=r$ gives $r=t$.
\end{proof}

Symmetry follows by joining $s=s$ and $s=t$. To obtain transitivity
from $s=t$ and $t=u$, reverse the first equation and join the two
equations with common left side $t$.

Both computation equations are derived:
\[
 I(0,a;u.h)=a,
 \qquad I(Sc,a;u.h)=h[I(c,a;u.h)/u].             \tag{I0, IS}
\]
Use the appropriate Eval rule on reflexivities of its two sources.

\begin{lemma}[One-term transport]\label{lem:transport}
For any fixed term template $h(v)$, from $s=t$ and $h(s)=r$ derive
$r=h(t)$. Consequently all numerical congruence and equational
replacement rules are available.
\end{lemma}
\begin{proof}
Let $J=I(0,s;v.h(v))$ and $K=I(S0,s;v.h(v))$.
I0 gives $J=s$, hence $J=t$ by transitivity. Eval$_S$ on $J=s$ and
$K=K$ gives $K=h(s)$, hence $K=r$. Eval$_S$ on $J=t$ and $K=r$
now gives $r=h(t)$. Choosing the second premise as a reflexivity and
using symmetry yields ordinary congruence. Transport in both sides of
an equation is two applications with the corresponding templates.
\end{proof}

\section{Recovering cases before induction}

\begin{lemma}[Provable mirror bases]\label{lem:base}
If $\Mir(B)$, then $\Gamma\vdash B(t,0)$ for every context and witness.
\end{lemma}
\begin{proof}
The instance is $a(t,0)=I(0,a(t,0);u.h(t,0,u))$. This is the symmetry
of I0. No induction, coverage, or logical macro is used.
\end{proof}

\begin{lemma}[All case introductions]\label{lem:allcases}
For every raw matrix $B$, the following rule is admissible:
\[
 \frac{\Gamma,B(t,0)\vdash B(t,Sz)}
      {\Gamma\vdash D_{x,y}B}.
\]
\end{lemma}
\begin{proof}
If the matrix is not mirrored, this is primitive DI. If it is mirrored,
cut the proof from Lemma~\ref{lem:base} into the supplied branch, obtaining
$B(t,Sz)$ under $\Gamma$ alone. Weaken by $B(t,z)$ and use progressive
DI. The step variable remains fresh because the removed base contains
no $z$.
\end{proof}

\begin{lemma}[Case-closure introduction and numerical cases]\label{lem:ncases}
For every native formula $F$, both rules are admissible:
\[
 \frac{\Gamma,F(0)\vdash F(Sz)}{\Gamma\vdash Q_yF(y)},
 \qquad
 \frac{\Gamma\vdash F(0)\qquad\Gamma\vdash F(Sz)}
      {\Gamma\vdash F(v)}.
\]
In the second rule $z$ is fresh for the context, passive parameters,
and requested target $v$.
\end{lemma}
\begin{proof}
Introduce $D_{w,y}F(y)$ by the derived case rule with an unused witness
coordinate $w$ and witness $0$. DE exposes $Q_yF(y)$, independent of
the fresh witness, so Hyp returns that closure and DE closes. This is
the first rule. For the second, weaken the successor premise by $F(0)$,
apply the first rule, and apply QE to the supplied base at $v$.
\end{proof}

The case rule is available for arbitrary compound formulas. Its proof
uses no equation induction. In particular, the numerical-case arguments
in the next section do not assume the induction they are used to recover.

\section{Recovering unrestricted equation induction}\label{sec:induction}

Write $C(c,a,b)=I(c,a;u.b)$, with $u$ unused. Computation gives
$C(0,a,b)=a$ and $C(Sv,a,b)=b$. Let
\[
 E(y)\equiv s(y)=r(y),\qquad
 G(y)\equiv s(y)=C(y,s(y),r(y)).                      \tag{Lift}
\]
The latter is a mirrored equation with one iterator. It has a provable
zero instance and its successor instance is interderivable with that
of $E$, using only computation and equality.

\begin{lemma}[Equational transfer]\label{lem:transfer}
For every target $v$, from $E(0)$ and $G(v)$ derive $E(v)$ without
equation induction.
\end{lemma}
\begin{proof}
Numerical cases establish the following unconditional equations:
\begin{align}
 C(v,s(0),s(v))&=s(v),                               \label{eq:Ls}\\
 C(v,r(0),C(v,s(v),r(v)))&=r(v).                     \label{eq:Lr}
\end{align}
At zero both left sides compute to the displayed zero seed. At a
successor the first equation uses one constant-update computation and
the second uses two. Apply Lemma~\ref{lem:ncases} to these equations
\emph{before} retaining any assumption depending on $v$. Then weaken
the proved equations into the desired context. This order is necessary
for the eigenvariable condition.

Under $G(v)$, congruence for the template $C(v,s(0),-)$ gives
\[
 C(v,s(0),s(v))=C(v,s(0),C(v,s(v),r(v))).
\]
The retained base $s(0)=r(0)$ changes the outer seed on the right.
Using \eqref{eq:Ls} and \eqref{eq:Lr}, the complete chain is
\begin{align*}
 s(v)
 &=C(v,s(0),s(v))\\
 &=C(v,s(0),C(v,s(v),r(v)))\\
 &=C(v,r(0),C(v,s(v),r(v)))\\
 &=r(v).
\end{align*}
Each step is a previously available equational inference.
\end{proof}

\begin{theorem}[Contextual equation induction]\label{thm:induction}
For every literal equation $E$ and arbitrary retained context $\Gamma$,
\[
 \frac{\Gamma\vdash E(0)\qquad\Gamma,E(z)\vdash E(Sz)}
      {\Gamma\vdash E(t)}
\]
is admissible, with the usual fresh $z$.
\end{theorem}
\begin{proof}
Use the lift $G$ above. Under $\Gamma,G(z)$, the base proof and
Lemma~\ref{lem:transfer} give $E(z)$. Cut this into the supplied step
derivation to get $E(Sz)$. The constant-update computation converts this
to $G(Sz)$. Thus the progressive premise for $G$ has been obtained.

With an unused existential coordinate $w$ and witness $0$, primitive
DI gives $D_{w,y}G(y)$. DE exposes $Q_yG(y)$. The theorem $G(0)$
and QE at $t$ yield $G(t)$; the retained base and
Lemma~\ref{lem:transfer} yield $E(t)$. Close DE. No substitution into
the supplied step derivation was needed, apart from injective fresh
renaming of its eigenvariable.
\end{proof}

\begin{corollary}[Native equation-progress rule]\label{cor:qprogress}
From $\Gamma,E(z)\vdash E(Sz)$ derive $\Gamma\vdash Q_yE(y)$.
\end{corollary}
\begin{proof}
Use the first rule of Lemma~\ref{lem:ncases}. In its branch retain
$E(0)$ and apply Theorem~\ref{thm:induction} at target $Sv$ to obtain
the required $E(Sv)$. The step eigenvariable is renamed fresh from
this outer $v$ and all retained data. Then discharge $E(0)$ by the
case-closure rule.
\end{proof}

\section{Exact sequent equivalence and numerical substitution}

\subsection{Proof of the main equivalence}
For $\Mcal\to\Ncal$, Hyp, Refl, Eval, QE, and DE have the existing
native derivations. A nonmirrored DI is the native case introduction.
A mirrored DI has a literal equation matrix. The native equation-progress
rule gives $Q_yB(t,y)$; under $B(t,0)$, QE supplies $B(t,Sz)$, and
native case introduction with witness $t$ gives $D_{x,y}B$. This
translates every primitive rule without changing the concluding sequent.

For $\Ncal\to\Mcal$, derive I0, IS, and one-term transport as in
Section~3. Hyp, QE, and DE are primitive. Replace native case introduction
by Lemma~\ref{lem:allcases} and native equation-progress introduction by
Corollary~\ref{cor:qprogress}. Induction on the native derivation completes
the transformation. Together these prove the first assertion of
Theorem~\ref{thm:main}. The existing equality of \texttt{SPrf} and
\texttt{NPrf} gives exact equality with the six-rule baseline as well.

\subsection{Substitution may change the mode}
\begin{theorem}[Direct substitution admissibility]\label{thm:subst}
If $\Gamma\vdash_{\Mcal}A$ and $\sigma$ is any capture-free substitution
of terms for passive numerical variables, then
$\Gamma\sigma\vdash_{\Mcal}A\sigma$.
\end{theorem}
\begin{proof}
The mirror test is preserved by substitution: the repeated initializer
and left side receive the same substitution, while the distinguished
bound counter remains $y$. It need not be reflected. For example,
$p=I(y,q;u.0)$ becomes mirrored after both $p,q$ are replaced by $0$.

Induct on the derivation. Rename each eigenvariable fresh from the finite
union of free variables in the relevant substitution images. All rules
other than DI commute with the corresponding lifted substitutions. If
DI originally used progress, preservation gives the same mode after
substitution, so apply the induction hypothesis and primitive DI. If
DI originally used cases, the induction hypothesis gives a substituted
case premise; apply the already derived unrestricted case introduction
of Lemma~\ref{lem:allcases}. That rule repairs a newly mirrored case by
cutting its provable zero base and weakening to progress.
\end{proof}

This proof is internal to the new calculus; it does not translate to
$\Ncal$ and back. The repair principle is inherited from
\texttt{reflexive\_base} \cite{reflexive}. Unlike the current diagonal
selector, the mirror selector is not preserved \emph{and reflected} by
arbitrary substitutions. The theorem claims admissibility with an
explicit proof repair, not that the primitive DI instance always keeps
the same mode.

% Review contribution: ordinary realization and inherited bi-interpretation.
% Not a new arithmetic-bootstrap proof. Citation keys may be adapted.
\section{Ordinary realization and the precise iPRA equivalence}

The new result changes one introduction selector.  Its arithmetic bridge is
the existing ConciseIPRA bridge, and is transferred by equality of the
\emph{entire contextual consequence relation}.  This distinction prevents
agreement on the standard natural numbers from being mistaken for
bi-interpretability.

\subsection{A completely specified ordinary realization}
Elaborate iterator syntax into a staged first-order signature.  Start with
$0,S$.  For each fixed earlier-stage term $h(u,\vec p)$ introduce a symbol
$i_h(n,a,\vec p)$.  Its defining axioms are the universal closures of
\[
 i_h(0,a,\vec p)=a,\qquad
 i_h(Sn,a,\vec p)=h(i_h(n,a,\vec p),\vec p).
\]
Only finite descriptions are used.  There is no variable ranging over
function descriptions and no universally total evaluator.

Let $T_I$ have ordinary intuitionistic first-order logic with equality,
the displayed computation schemas, the two axioms
\[
 \forall x\bigl(x=0\lor\exists y\,x=Sy\bigr),\qquad
 (0=S0)\to\bot,
\]
and contextual literal-equation induction
\[
 \frac{\Gamma\vdash E(0)\qquad\Gamma,E(z)\vdash E(Sz)}
      {\Gamma\vdash E(t)}.
\]
Here $E$ is a literal equation and $z$ is fresh for the retained context,
the target, and the passive parameters.  No induction on a quantified
formula is primitive.

For completeness, the ordinary logic consists of reusable hypotheses;
$\land$ introduction and both projections; both $\lor$ introductions and
two-branch elimination; implication introduction discharging one
assumption and modus ponens; universal introduction at a fresh
eigenparameter and universal specialization at any term; existential
introduction at a term and existential elimination at a fresh
eigenparameter; falsity elimination; equality reflexivity; and equality
replacement in any formula.  Context weakening, exchange, contraction,
capture-free substitution, and cut are admissible.  Freshness for
universal introduction concerns the retained context; freshness for
existential elimination concerns both the retained context and its
conclusion.  These are all the logical and nonlogical primitives of
$T_I$.

The ordinary reading of a native formula is
\[
 (s=t)^\circ=(s=t),\qquad
 (Q_yB)^\circ=B^\circ(0)\to\forall y\,B^\circ(y),
\]
\[
 (D_{x,y}B)^\circ=
 \exists x\bigl(B^\circ(x,0)\to\forall y\,B^\circ(x,y)\bigr).
\]
The interpretation domain is $x=x$ and equality is identity.  The
designated native contradiction $0=S0$ interprets as an actual
contradiction because the separation axiom is explicit.  Omitting that
axiom would leave the ordinary positive arithmetic axioms compatible
with a singleton structure.

The new mirrored introduction is sound in $T_I$: in its progressive
branch its matrix is a literal equation, so contextual equation
induction applies under the assumed base.  In its case branch numerical
coverage combines the zero case and all successor cases.  The other
five rule families have their displayed equality or ordinary logical
readings.

\subsection{Explicit inverse logical macros}\label{sec:choice-introduction}
One can first erase auxiliary $Q$ occurrences by replacing
$Q_yB(y)$ with $D_{u,y}B(y)$, with fresh unused $u$.  These formulas are
interderivable: case introduction gives the forward implication, and
paired elimination returns the opened closure in the reverse direction.
Congruence under $Q$ and $D$ then gives this replacement inside every
native formula.  Thus the following macros need only the Eq/$D$
fragment, without restricting the public language.

Write $C(v,a,b)=I(v,a;u.b)$, with $u$ unused, and $1=S0$.  On Eq/$D$
formulas define the gate
\[
 G_v(s=t):=\bigl(C(v,0,s)=C(v,0,t)\bigr),\qquad
 G_v(D_{x,y}A):=D_{x,y}G_vA,
\]
renaming all binders away from $v$.  Then $G_0A$ is a theorem,
$G_{Sv}A\dashv\vdash A$, and
$A\dashv\vdash D_{u,y}G_yA$ for fresh unused $u,y$.

For two equations set
\[
 M_v(s=t,r=w):=\bigl(C(v,s,r)=C(v,t,w)\bigr).
\]
Extend by aligned fresh paired binders:
\[
\begin{aligned}
 M_v(D_{x,y}A,D_{x,y}B)&:=D_{x,y}M_v(A,B),\\
 M_v(E,D_{x,y}B)&:=D_{x,y}M_v(G_yE,B),\\
 M_v(D_{x,y}A,E)&:=D_{x,y}M_v(A,G_yE),
\end{aligned}
\]
where $E$ is an equation.  Each recursive call reduces maximum binder
depth.  The endpoint equivalences are
$M_0(A,B)\dashv\vdash A$ and $M_{Sv}(A,B)\dashv\vdash B$.

Put $\delta(v)=C(v,1,0)$ and $p(z)=I(z,0;u.\delta(u))$.  Only
$p(1)=1$ and $p(2)=0$ are needed.  The full inverse logical translation is
\[
\begin{aligned}
 A\Rightarrow_*B&:=D_{u,z}M_z(A,B),\\
 A\wedge_*B&:=D_{u,z}G_zM_{p(z)}(A,B),\\
 \exists_*x\,A(x)&:=D_{x,y}G_yA(x),\\
 A\vee_*B&:=\exists_*x\,M_x(A,B),\\
 \forall_*x\,A(x)&:=A(0)\wedge_*D_{u,z}G_zA(z),\\
 \bot_*&:=(0=1).
\end{aligned}
\]
All newly bound variables are fresh, and $u$ is unused.

Case introduction and the selector endpoints prove implication
introduction; elimination opens its clause and specializes at $1$.
The gate proves existential introduction and elimination.  Disjunction
uses witnesses $0$ and $1$, while its elimination uses the derived
numerical case rule.  Conjunction projections specialize at $1$ and
$2$; introduction uses numerical cases for the selector.  The first
universal conjunct supplies zero and the second supplies successors,
so coverage proves universal elimination.  Finally $0=1$ gives
$s=C(0,s,t)=C(1,s,t)=t$ for arbitrary terms; case introduction then
extends explosion to every formula.  These are the inherited logical
macro proofs, applicable because all case introductions and all
contextual equation inductions have been recovered in the new calculus.

The essential native formula composites are internally interderivable:
\[
 D_{x,y}B\dashv\vdash
 \exists_*x\bigl(B(x,0)\Rightarrow_*\forall_*y\,B(x,y)\bigr),
\]
\[
 Q_yB\dashv\vdash B(0)\Rightarrow_*\forall_*y\,B(y).
\]
For the first equivalence, paired elimination exposes the conditional
universal at a fresh witness; the converse uses paired case introduction
at the existential witness.  The second follows by closure elimination
and derived case-closure introduction.  External structural induction
on finite formulas proves both complete logical round trips.

\subsection{The inherited arithmetic bridge, with both composites}
The ConciseIPRA arithmetic development proves the following lemmas~\cite{bridge,readme}
\emph{internally in $T_I$}, in this order: a scalar predecessor
$P0=0$, $P(Sn)=n$; a term reflecting actual equality; a characteristic
term for every quantifier-free formula; pairing with two projections;
and a compiler for every fixed primitive-recursive description.  These
are inherited results, not new consequences asserted solely from
finite tests.  The source proofs are the reference files
\texttt{scalar\_predecessor.tex}, \texttt{equality\_reflection.tex}, and
\texttt{primitive\_recursion.tex}, with their corresponding Lean
modules named in the pinned repository README.

Here is the exact compiler step.  Suppose pairing $\pi$ and projections
$L,H$ have already been established, and a target PR description is
\[
 f(0,\vec p)=g(\vec p),\qquad
 f(Sn,\vec p)=h(n,f(n,\vec p),\vec p).
\]
For translations $g^*,h^*$ of its earlier descriptions put
\[
 V(n,\vec p)=I\bigl(n,\pi(0,g^*(\vec p));
 u.\pi(SL(u),h^*(L(u),H(u),\vec p))\bigr),
 \qquad f^*(n,\vec p)=H(V(n,\vec p)).
\]
Equation induction proves $L(V(n,\vec p))=n$.  The two projection
identities then establish both required defining equations for $f^*$.
The compiler recurses externally on a fixed finite PR description.

Conversely, iPRA interprets each iterator by primitive recursion whose
update ignores the old counter.  It proves the computation equations,
coverage, separation, and every literal-equation induction instance.
In the other direction, the predecessor gives successor injectivity,
the zero test gives successor separation, and the characteristic terms
reduce quantifier-free induction to literal-equation induction.

For the numerical round trips, two recursive operations with the same
seed and provably equal updates are internally equal: induction on the
literal equation between their outputs proves this at every counter.
External induction on finite descriptions therefore shows that every
operation is provably equal to its double translation.  Both domains
are whole and both interpreted equalities are actual equality, so
$x=y$ is an internally definable isomorphism from each structure to its
double interpretation.  Formula induction, using the logical composites
above, gives both formula round trips.  Effective proof translations
preserve contradiction in both directions and hence give
equiconsistency; they make no assertion that either theory proves its
own consistency.

Consequently the new calculus is a faithful presentation of an ordinary
theory bi-interpretable and equiconsistent with full first-order iPRA.
The proof of this conclusion uses the new exact sequent equivalence and
the inherited arithmetic bridge.  It does not rely on a standard-model
test, a quotient equality, or an interpretation restricted to standard
numerals.


\section{Checker, exact comparison, and verification}\label{sec:implementation}

\subsection{Public certificate contract}
The standalone checker uses Python 3.10 or later and only the standard
\texttt{json} and \texttt{sys} modules. The measured runtime is CPython
3.12.14. Input is one JSON object with exactly four fields:
\begin{lstlisting}
{"variables": n, "context": [A0, A1, ...], "goal": A, "proof": p}
\end{lstlisting}
Here $n$ is an exact nonnegative integer, not a Boolean. Terms use
\texttt{["v",i]}, \texttt{["z"]}, \texttt{["s",t]}, and
\texttt{["i",n,a,h]}; formulas use \texttt{["=",s,t]},
\texttt{["q",B]}, and \texttt{["d",B]}. Binding is de Bruijn:
the distinguished $y$ is index $0$ and $x$ index $1$ in a $D$ matrix.
The iterator increments binding depth only in its update body.

The proof tags are \texttt{hyp}, \texttt{refl}, \texttt{eval},
\texttt{qe}, \texttt{de}, and \texttt{di}. Tag fields and every supplied
term and formula are checked at the appropriate scope; malformed active
syntax, Boolean variable indices, invalid references, and escaping
eigenvariables are rejected. Each recursive inference computes its
conclusion and the top-level goal is checked against that result.
The only successful public output is \texttt{ACCEPT} followed by a
newline and exit status zero. A rejected input exits unsuccessfully.

The checker is a direct implementation of Section~\ref{sec:spec}.
Validation of terms/formulas and a structural induction on successful
inference calls establish that accepted certificates are derivations of
the specified calculus. Conversely, every finite derivation has the
corresponding tree certificate. The completeness statement, like that
of the repository checkers, assumes unbounded available resources;
a fixed Python recursion or memory limit may reject a very large valid
certificate. Such a failure is never acceptance.

\subsection{Matched measurements}
The comparator implements the current \texttt{SPrf} calculus in the
same typed Python framework. It is a supplied matched implementation,
not a claim that this exact Python file already appears upstream.
Both files have the same parser, term and formula validators,
capture-avoiding substitution, context interface, and all other rule
branches. The only executable-source difference is the DI guard.

\begin{center}
\begin{tabular}{lrrr}
\toprule
Complete normalized module & Bytes & Tokens & AST nodes\\
\midrule
Current diagonal selector & 3125 & 1022 & 1091\\
Proposed mirror selector  & 3109 & 1012 & 1082\\
Decrease                 &   16 &   10 &    9\\
\bottomrule
\end{tabular}
\end{center}

The measurement removes genuine docstrings, applies \texttt{ast.unparse},
and appends one line feed. Bytes mean normalized UTF-8 bytes. Significant
tokens exclude the end marker, encoding, comments, newlines, indentation,
and dedentation. AST nodes are counted after docstring removal. The
archive contains the complete source, measurement code, normalized
copies, and a one-line executable diff.

This is a 0.51\% byte reduction against the matched six-rule control.
It is not a reduction below the repository's 516-byte compact codec,
nor a new minimum among all of its public-goal or fully validated
interfaces. Those smaller programs use different interfaces and
encodings. A six-rule presentation and a shorter matched recognizer
are the supported simplicity claims.

The tradeoff is also concrete: the diagonal guard inspects fixed
constructor/variable positions, whereas the mirror guard may compare
two large term trees. That comparison is linear in their size in the
worst case. The source reduction does not imply uniformly faster
checking or shorter proofs for every theorem.

\subsection{Evidence and limits}
The executable evidence includes the following checks; their complete
results are recorded in the accompanying JSON ledgers.
\begin{center}
\begin{tabular}{lr}
\toprule
Check & Executions or cases\\
\midrule
Direct fixtures, with expected outcomes for both kernels & 23\\
Direct API checks across both kernels & 46\\
Standalone checks, ordinary Python and \texttt{-O} & 92\\
Independent replay in the fetched native checker & 18\\
Native standalone replays, ordinary Python and \texttt{-O} & 36\\
Numerical-substitution replays & 72\\
\bottomrule
\end{tabular}
\end{center}
These counts describe repeated executions across interfaces and
implementations, not that many independent mathematical phenomena.
The fixtures exercise both computations; symmetry, transitivity, and
congruence with repeated occurrences and a binder; a progressive proof
that consumes its induction hypothesis; closure specialization at
$0,1,2$; compound introduction; a nonzero witness occurring in both
mirrored terms; nested iterator binders; and scope and malformed-input rejection.
An intentional mutant that tests the mirror shape after witness
substitution is distinguished by the before-witness boundary fixture.

The untrusted transformation code implements derived case introduction,
simultaneous numerical and hypothesis substitution, and translation
from the mirror calculus to the unchanged fetched native checker.
In particular, the source matrix
\[
 B(x,y)\equiv p=I(y,q;u.u)
\]
is unselected when $p,q$ are distinct variables, but becomes selected
after both are substituted by the same variable. Its source branch
uses its zero-base hypothesis. The literally substituted unrepaired
proof is rejected and the repaired proof is accepted. A second test
repairs a mirrored case introduction whose old branch also consumes
its zero base. Both deliberate unrepaired controls are rejected.

The reverse mirror-to-native compiler is generic. A generic executable
compiler for every native equation-progress proof in the other
direction is not included; that direction is established by the
written construction in Section~\ref{sec:induction}.


The general equivalence and interpretation claims are mathematical
proofs in this note. Executed certificates support the implementation
and selected proof transformations; finite testing cannot establish
the universal metatheorems. The existing Lean source endpoints were
read at the pinned revision. Their project was not rebuilt here, and
there is no new Lean or Rocq formalization of the mirror calculus or
of the Python checker. The included Lean files are reference excerpts
with their original imports, not a standalone formalization project.

\section{Further research questions}
\begin{enumerate}
\item Can the same selector be integrated into the repository's smallest
complete certificate interface with a net source reduction after all
validation and scope checks are counted? The present comparison does
not answer that question.
\item What is the smallest family of syntactic base recognizers that
supports an equational decoder for every equation? The mirror and
diagonal shapes give two concrete instances, but neither establishes
a lower bound.
\item Can the mirror comparison be made constant-time by syntax sharing
without moving a larger validation obligation into the trusted input
reader? Any complexity comparison must count the representation and
its integrity checks.
\item Can the equation-induction translator be formalized directly in
Lean with fewer proof dependencies than the current diagonal
construction? The two unconditional selector equations are a useful
target for such a formalization.
\item Can one prove a substantive lower bound for a precisely specified
class of complete proof checkers? Absolute claims over unrestricted
programming languages or encodings are not supported by the present
experiments.
\end{enumerate}

\appendix
\section{Complete standalone checker}\label{app:checker}
The following is the complete selected acceptance module. No proof
generator, source translator, reference theorem, or measurement helper
is called by it.
\begin{lstlisting}[language=Python]
"""Six-rule mirror-guard proposal, with raw terms and two-case Eval."""
import json
import sys


def require(condition):
    if not condition:
        raise ValueError('Invalid certificate')


def term(a, n):
    match a:
        case ['v', i]:
            require(type(i) is int and 0 <= i < n)
        case ['z']:
            pass
        case ['s', b]:
            term(b, n)
        case ['i', c, b, h]:
            term(c, n)
            term(b, n)
            term(h, n + 1)
        case _:
            raise ValueError('Invalid term')


def formula(a, n):
    match a:
        case ['=', s, t]:
            term(s, n)
            term(t, n)
        case ['q', b]:
            formula(b, n + 1)
        case ['d', b]:
            formula(b, n + 2)
        case _:
            raise ValueError('Invalid formula')


def walk(a, f, d=0):
    if a[0] == 'v':
        return f(a[1], d)
    return [a[0]] + [walk(b, f, d + ({'i': [0, 0, 1], 'q': [1], 'd': [2]}.get(a[0], [0] * len(a)))[j])
                     for j, b in enumerate(a[1:])]


def shift(a, k=1):
    return walk(a, lambda i, d: ['v', i + k if i >= d else i])


def sub(a, t, drop=1):
    return walk(a, lambda i, d: ['v', i] if i < d else shift(t, d) if i == d else ['v', i - drop])


def infer(p, g, n):
    match p:
        case ['hyp', i]:
            require(type(i) is int and 0 <= i < len(g))
            return g[i]
        case ['refl', a]:
            term(a, n)
            return ['=', a, a]
        case ['eval', p, q]:
            e, f = infer(p, g, n), infer(q, g, n)
            require(e[0] == f[0] == '=' and f[1][0] == 'i')
            s, t = e[1:]
            _, c, a, h = f[1]
            if c == ['z']:
                require(s == a)
                return ['=', f[2], t]
            require(c[0] == 's' and s == ['i', c[1], a, h])
            return ['=', f[2], sub(h, t)]
        case ['qe', t, p, q]:
            term(t, n)
            a = infer(p, g, n)
            require(a[0] == 'q')
            require(infer(q, g, n) == sub(a[1], ['z']))
            return sub(a[1], t)
        case ['de', a, p, q]:
            formula(a, n)
            b = infer(p, g, n)
            require(b[0] == 'd')
            require(infer(q, [['q', b[1]]] + [shift(x) for x in g], n + 1) == shift(a))
            return a
        case ['di', b, t, p]:
            formula(b, n + 2)
            term(t, n)
            a = sub(['q', b], t)[1]
            h = a if b[0] == '=' and b[2][:3] == ['i', ['v', 0], b[1]] else shift(sub(a, ['z']))
            require(infer(p, [h] + [shift(x) for x in g], n + 1) == sub(a, ['s', ['v', 0]], 0))
            return ['d', b]
        case _:
            raise ValueError('Invalid proof')


def check(c):
    require(type(c) is dict and set(c) == {'variables', 'context', 'goal', 'proof'})
    n, g, a, p = (c[k] for k in ['variables', 'context', 'goal', 'proof'])
    require(type(n) is int and n >= 0 and type(g) is list)
    for b in g:
        formula(b, n)
    formula(a, n)
    require(infer(p, g, n) == a)


if __name__ == '__main__':
    try:
        check(json.load(sys.stdin))
    except (ValueError, TypeError, IndexError, KeyError, RecursionError):
        sys.exit(1)
    print('ACCEPT')
\end{lstlisting}

\clearpage
\section*{Inherited arithmetic appendices}
\addcontentsline{toc}{section}{Inherited arithmetic appendices}
The following three sections reproduce the pinned ConciseIPRA
predecessor, equality-reflection, and primitive-recursion arguments
\cite{erasure,forward,progress}. They are inherited foundations rather
than new reductions. They make the dependency of the iPRA bridge
reviewable without relying on standard-model agreement. The original
source files and their license are also included in the archive.
% Adapted with attribution from the pinned unparenthesized one-hole archive.
\section{Internal scalar predecessor bootstrap}
\label{sec:hybrid-bootstrap}

This scalar bootstrap is inherited from the supplied scalar development~\cite{erasure} and is reproduced to keep the scalar checker self-contained. It is not a new contribution of this paper. The proof uses a factor map fixing odd states and a saturated-iteration construction of predecessor.

The available arithmetic terms are $0$, $S$, and
\[
 I(n,a;u.h),\qquad
 I(0,a;u.h)=a,\qquad
 I(Sn,a;u.h)=h[I(n,a;u.h)/u].
\]
The update binds only $u$.  All other variables in it are fixed numerical
parameters.  We use equality reasoning, the already reconstructed
zero/successor cases, and induction on a single literal equation.
Every induction below keeps its other parameters fixed.  Predecessor,
successor injectivity, pairing, and induction on arbitrary
quantifier-free formulas are not assumptions of the argument.

\subsection{Elementary terms and two finite cycles}

Write $1=S0$ and use larger numerals as successor abbreviations.  Define
\[
 C(n,a,b)=I(n,a;u.b),\qquad
 Z(n)=C(n,1,0),\qquad
 a+b=I(b,a;u.Su).
\]
The update variable of $C$ is absent from its body.  Consequently
\[
 C(0,a,b)=a,\qquad C(Sn,a,b)=b.
\]
Numerical cases show that $Z(n)$ is $0$ or $1$.
The usual commutative-monoid equations for addition follow by literal
equation induction.  For example, $0+b=b$ and $Sa+b=S(a+b)$ are
proved by induction on $b$; these yield commutativity by a further
equation induction, and associativity follows by induction on the last
summand.  We use these equations to reorder finite sums.  The notation
$2a$ means $a+a$ and $3a$ means $a+a+a$; no multiplication operation
is needed for this bootstrap.

Put
\[
 N(n)=I(n,0;u.Z(u)),\qquad
 j(u)=Z(u)+2N(u),\qquad
 M(n)=I(n,0;u.j(u)).
\]
The range of $N$ is contained in $\{0,1\}$: its zero value is $0$,
and its successor value is a value of $Z$.  This uses numerical cases
on its counter, not induction on a disjunction.  Hence
\begin{equation}
 N(n)+Z(N(n))=1,\qquad
 Z(Z(N(n)))=N(n),\qquad N(n+2)=N(n).
 \label{eq:hybrid-parity}
\end{equation}
Here the last equation follows from two successor computations and
the preceding equation.  Literal-equation inductions give
\begin{equation}
 N(2n)=0,\qquad N(2n+1)=1,\qquad N(3n)=N(n).
 \label{eq:hybrid-parity-scaling}
\end{equation}
For the first equation the successor step adds two to the argument of
$N$; the second follows by one computation.  In the third, increasing
$n$ adds three to the left argument and therefore applies $Z$ once to
both sides of the induction equation.

Finite computations now give
\[
 j(0)=1,\qquad j(1)=2,\qquad j(2)=0.
\]
The proof below needs just these three values and
$M(Sn)=j(M(n))$; it does not assume a division algorithm or a theorem
characterizing $M$ as remainder modulo three.

\subsection{A ramp with a global scaling identity}

Define
\[
 G(n)=I(n,0;u.u+M(u)+1).
\]
The central invariant is
\begin{equation}
 M(G(n))=N(n).
 \label{eq:hybrid-ramp-invariant}
\end{equation}
Its zero case is immediate.  For its successor step, retain the
single induction equation and split the Boolean value $N(n)$.
If it is $0$, the next $G$ state is $G(n)+1$; its $M$ value is
$j(0)=1$, equal to $N(Sn)$.  If it is $1$, the next $G$ state
is $G(n)+2$; two computations of $M$ change its value from $1$
to $2$ and then to $0$, again equal to $N(Sn)$.  This proves
the invariant by literal-equation induction.

It follows that
\[
 G(Sn)=G(n)+N(n)+1,\qquad G(n+2)=G(n)+3.
\]
Equation induction therefore yields
\begin{equation}
 G(2n)=3n.
 \label{eq:hybrid-ramp-double}
\end{equation}
A separate literal-equation induction proves
\begin{equation}
 G(3n)=3G(n)+N(n).
 \label{eq:hybrid-ramp-triple}
\end{equation}
For completeness, its successor calculation is
\begin{align*}
 G(3Sn)
 &=G(3n)+N(3n)+4
   =3G(n)+2N(n)+4,\\
 3G(Sn)+N(Sn)
 &=3G(n)+3N(n)+3+Z(N(n)).
\end{align*}
The expressions agree in both Boolean cases for $N(n)$.
The base is $0=0$.  In particular,
\eqref{eq:hybrid-ramp-triple} was proved internally; it was not
inferred from an intended floor-function interpretation of $G$.

Set
\[
 H(u)=C(N(u),G(u),u).
\]
The crucial choice is to retain $u$ when its parity is $1$.
The preceding equations give
\begin{equation}
 H(2u)=3u,\qquad H(1)=1,\qquad H(3u)=3H(u).
 \label{eq:hybrid-factor-map}
\end{equation}
To check the last identity, use $N(3u)=N(u)$ and split the Boolean
$N(u)$.  At $0$, both sides are $3G(u)$ by
\eqref{eq:hybrid-ramp-triple}; at $1$, both are $3u$.
This argument does not require that an arbitrary number has already
been represented as $2v$ or $2v+1$.

\subsection{A comparator from scalar iteration}

For every fixed update $h$, equation induction on $n$, keeping $m,a$
and all parameters fixed, proves the semigroup identity
\[
 I(m+n,a;h)=I(n,I(m,a;h);h).
\]
The base is the zero computation and the step applies $h$ to the
induction equation. This identity will also be used in the later
pairing construction.


For a fixed update $h$, equation induction on $y$, with its initializer
$a$ and all parameters of $h$ fixed, proves the rotation identity
\begin{equation}
 I(Sy,a;h)=I(y,h(a);h).
 \label{eq:hybrid-rotation}
\end{equation}
Both sides have base $h(a)$, and their successor equations apply the
same $h$ to the previous equality.  Here and below $I(y,a;h)$ is
an abbreviation for one fixed finite scalar term, not quantification
over functions.

By \eqref{eq:hybrid-factor-map}, another equation induction gives
\begin{equation}
 I(y,3a;H)=3I(y,a;H).
 \label{eq:hybrid-iterator-scaling}
\end{equation}
Its base is $3a=3a$.  In the step the fixed-parameter induction
equation rewrites the new left side as
\[
 H(3I(y,a;H))=3H(I(y,a;H)).
\]
The initializer parameter stays $a$ throughout this induction.
Equation induction also gives $I(y,1;H)=1$, since $H(1)=1$.

Define
\[
 E(x)=I(x,1;u.2u),\qquad
 L(x,y)=N\bigl(I(y,E(x);H)\bigr).
\]
Then the following entire recursion specification is provable:
\begin{equation}
 L(0,y)=1,\qquad L(Sx,0)=0,\qquad L(Sx,Sy)=L(x,y).
 \label{eq:hybrid-comparator}
\end{equation}
The first equation uses $E(0)=1$ and the fixed-point identity for
$H$.  The second uses $E(Sx)=2E(x)$ and $N(2a)=0$.
For the third,
\begin{align*}
 L(Sx,Sy)
 &=N\bigl(I(Sy,2E(x);H)\bigr)\\
 &=N\bigl(I(y,H(2E(x));H)\bigr)\\
 &=N\bigl(I(y,3E(x);H)\bigr)\\
 &=N\bigl(3I(y,E(x);H)\bigr)\\
 &=N\bigl(I(y,E(x);H)\bigr)=L(x,y).
\end{align*}
The successive steps use rotation, $H(2a)=3a$, iterator scaling,
and $N(3a)=N(a)$.  Thus $L$ has the comparison equations at all
arguments internally.  Its standard-model meaning as the indicator
of $x\le y$ is motivation, not a premise of the following argument.

\subsection{Saturated iteration gives predecessor}

We now need only the three equations
\eqref{eq:hybrid-comparator}.  Define
\[
 B_m(u)=C(L(Su,m),u,Su),\qquad
 W(t,m)=I(t,0;u.B_m(u)).
\]
The subscript in $B_m$ records a fixed parameter; $B_m$ is an open
term abbreviation, not a function variable.

Numerical cases on $e$ prove the selector equation
\[
 C(e,Sa,Sb)=S C(e,a,b).
\]
Equivalently one may use equation induction on $e$, whose induction
hypothesis is unused.  Therefore
\begin{align}
 B_{Sm}(Su)
 &=C(L(SSu,Sm),Su,SSu)\notag\\
 &=C(L(Su,m),Su,SSu)=S B_m(u),
 \label{eq:hybrid-conjugacy}\\
 B_{Sm}(0)&=C(L(S0,Sm),0,S0)=C(L(0,m),0,S0)=S0.
 \label{eq:hybrid-conjugacy-base}
\end{align}

Literal-equation induction on $t$, with $m$ fixed, proves
\begin{equation}
 W(St,Sm)=S W(t,m).
 \label{eq:hybrid-saturated-shift}
\end{equation}
The base is \eqref{eq:hybrid-conjugacy-base}.  The successor step is
\begin{align*}
 W(SSt,Sm)
 &=B_{Sm}(W(St,Sm))\\
 &=B_{Sm}(S W(t,m))\\
 &=S B_m(W(t,m))=S W(St,m).
\end{align*}
Only the second line uses the current induction hypothesis, at the
same fixed $m$.

Two subsequent equation inductions give
\begin{equation}
 L(Sn,n)=0,\qquad W(n,n)=n.
 \label{eq:hybrid-diagonals}
\end{equation}
The first uses the comparator recurrence.  The second uses
\eqref{eq:hybrid-saturated-shift} in its successor step.  The
substitution $t=m=n$ is made into that already proved universally
parameterized equation; no parameter of an active induction hypothesis
is changed.  Consequently
\begin{equation}
 W(Sn,n)=B_n(n)=C(L(Sn,n),n,Sn)=n.
 \label{eq:hybrid-extra-step}
\end{equation}

Finally define the finite scalar term
\begin{equation}
 P(m)=I\bigl(m,0;u.C(L(SSu,m),u,Su)\bigr).
 \label{eq:hybrid-predecessor-term}
\end{equation}
Its zero equation is immediate.  To justify its successor equation
without assuming a separate extensionality rule for iterator bodies,
put
\[
 V(t,m)=I\bigl(t,0;u.C(L(SSu,Sm),u,Su)\bigr).
\]
The comparator recurrence identifies this update pointwise with
$B_m(u)$.  Equation induction on $t$, with $m$ fixed, gives
$V(t,m)=W(t,m)$: their base values are $0$, and the successor step
uses the pointwise update equation and then the induction equation.
Now \eqref{eq:hybrid-extra-step} yields
\[
 \boxed{\qquad P(0)=0,\qquad P(Sm)=V(Sm,m)=W(Sm,m)=m.\qquad}
\]

\subsection{Dependency boundary}

All terms in this section are built from $0,S,I$ by finite substitution.
All inductions have a single equation as their invariant.  Numerical
cases and equality transport are the already reconstructed logical
resources.  Successor injectivity may now be obtained by applying $P$
to an assumed equation $Sx=Sy$.  Equality tests, reflection of
quantifier-free formulas, pairing, and the recovery of general
primitive recursion can then be developed in the established order.
They were not used in order to manufacture predecessor.

This argument is a written internal derivation in the stated theory.
The finite implementation checks do not constitute a proof-assistant
formalization of this entire metatheorem.

% Standalone section.  Requires only amsmath and amssymb in the parent document.
% No theorem environments, macros, or bibliography keys are assumed.
% All arithmetic claims in this section are proved here.

\section{From literal-equation induction to full intuitionistic PRA}
\label{sec:eic-completion}\label{app:arithmetic}

This section supplies an unconditional completion theorem. It is adapted
with attribution from the current forward-left draft~\cite{forward}. In the
scalar application used here, every function occurring in the proof is
already a scalar term: predecessor is supplied by
Appendix~\ref{sec:hybrid-bootstrap}, the selector is
$C(n,a,b)=I(n,a;u.b)$, addition is $x+y=I(y,x;u.Su)$, and
truncated subtraction is $d(x,y)=I(y,x;u.Pu)$. All other functions below
are finite compositions of these. Thus this proof does not presuppose
that arbitrary primitive recursion has already been compiled. The QFI
conclusion first applies to the scalar language; the general PR compiler
comes afterward.  Literal-equation
induction, together with the primitive-recursive defining equations,
predecessor coverage, and separation of $0$ and $1$, proves induction for every
quantifier-free formula.  In particular, the proof does not assume a faithful
numerical equality test, stability or decidability of equality, successor
uniqueness, or any commutation law for an iterator.  We construct the equality
test and prove its reflection property for arbitrary elements of the
numerical domain.

\subsection{The exact hypotheses}

Work in ordinary intuitionistic first-order logic with equality and genuine
falsity $\bot$.  The language has one numerical sort and a symbol for each
fixed, finite primitive-recursive description.  Distinct descriptions are
not identified merely because they compute the same standard function.
Let $\mathsf{PRE}$ contain the defining equations for zero, successor,
projections, composition, and primitive recursion.  In particular, choose
descriptions $P$ and $C$ with equations
\begin{align}
 P(0)&=0, & P(Sx)&=x, \label{eq:eic-predecessor}\\
 C(0,a,b)&=a, & C(Sx,a,b)&=b. \label{eq:eic-selector}
\end{align}
Write $1=S0$.  In addition to $\mathsf{PRE}$, assume
\begin{align}
 \mathsf{Cov}:&\qquad
       \forall x\,\bigl(x=0\lor x=S(Px)\bigr),
       \label{eq:eic-coverage}\\
 \mathsf{Sep}:&\qquad 0=1\longrightarrow\bot,
       \label{eq:eic-separation}
\end{align}
and induction for each \emph{literal equation} $E(n,\vec p)$:
\begin{equation}
 \forall\vec p\,
 \left[
   \left(E(0,\vec p)\land
       \forall n\bigl(E(n,\vec p)\longrightarrow E(Sn,\vec p)\bigr)
   \right)
   \longrightarrow\forall n\,E(n,\vec p)
 \right].
 \label{eq:eic-induction}
\end{equation}
Here the schematic expression $E$ is a single equality of numerical
terms; the displayed universal quantifiers range only over numbers.
Every description and every instance of the schema is finite external
syntax.

The axiom formulation \eqref{eq:eic-induction} and the usual contextual
induction rule are interchangeable.  For example, if
\[
 \Gamma\vdash E(0,\vec p),
 \qquad
 \Gamma,E(z,\vec p)\vdash E(Sz,\vec p),
\]
where $z$ is absent from $\Gamma$ and the passive parameters, discharge the
second premise's equation and universally introduce $z$.  The axiom then
gives $\Gamma\vdash\forall n E(n,\vec p)$, which can be instantiated at any
target term.  Conversely, assume the antecedent of the axiom and apply the
contextual rule with fresh variables.  This uses the ordinary intuitionistic
logical rules and permits arbitrary, including quantified, formulas in
$\Gamma$.

\subsection{Constructor separation and two elementary tests}

Define the sign description and the countdown description by
\begin{equation}
 \operatorname{sg}(u)=C(u,0,1),
 \qquad
 K(u)=C(u,1,Pu).
 \label{eq:eic-sign-countdown}
\end{equation}
Their defining equations imply
\begin{align}
 \operatorname{sg}(0)&=0,
 &\operatorname{sg}(Sx)&=1,\label{eq:eic-sign-computation}\\
 K(0)&=1,
 &K(Sx)&=x.\label{eq:eic-countdown-computation}
\end{align}

\paragraph{Constructor laws.}
Successor is injective: applying $P$ to $Sx=Sy$ gives $x=y$ by
\eqref{eq:eic-predecessor}.  Successor never equals zero.  Indeed,
$Sx=0$ gives $1=0$ by applying $\operatorname{sg}$ and using
\eqref{eq:eic-sign-computation}; symmetry and
\eqref{eq:eic-separation} give $\bot$.  Consequently $1\ne0$, where
$A\ne B$ abbreviates $A=B\to\bot$.  None of these arguments uses
induction.

\paragraph{The sign test.}
For every $u$,
\begin{align}
 &\operatorname{sg}(u)=0\lor\operatorname{sg}(u)=1,
       \label{eq:eic-sign-range}\\
 &\operatorname{sg}(u)=0\quad\longleftrightarrow\quad u=0,
       \label{eq:eic-sign-zero}\\
 &\operatorname{sg}(u)=1\quad\longleftrightarrow\quad u\ne0.
       \label{eq:eic-sign-one}
\end{align}
To prove the range assertion, split \eqref{eq:eic-coverage}.  The branch
$u=0$ gives sign $0$, and the branch $u=S(Pu)$ gives sign $1$.
The implication $u=0\to\operatorname{sg}(u)=0$ is computation.
For its converse, the successor branch of coverage would make
$\operatorname{sg}(u)=0$ imply $1=0$, a contradiction; in that branch
explosion gives $u=0$ as well.  Similarly, if
$\operatorname{sg}(u)=1$ and $u=0$, computation gives $0=1$.
Conversely, $u\ne0$ excludes the zero branch of coverage, and the
successor branch computes the sign as $1$.

Coverage and \eqref{eq:eic-selector} also give, for all $u,v$,
\begin{equation}
 C(u,v,v)=v.
 \label{eq:eic-selector-idempotence}
\end{equation}
Both coverage branches compute to $v$, so this identity again needs no
induction.

\paragraph{The countdown test and the notation $W$.}
For every $u$,
\begin{equation}
 K(u)=0\quad\longleftrightarrow\quad u=1.
 \label{eq:eic-countdown-zero}
\end{equation}
If $u=1=S0$, \eqref{eq:eic-countdown-computation} gives $K(u)=0$.
Conversely, assume $K(u)=0$ and split coverage.  If $u=0$, then
$K(u)=1$, contradicting separation, and explosion gives $u=1$.
If $u=S(Pu)$, computation gives $K(u)=Pu$; hence $Pu=0$ and
$u=S0=1$.

If a presentation uses a description named $W$ with defining equations
\[
 W(0)=1,\qquad W(Sx)=x,
\]
then it proves $K(u)=W(u)$ for every $u$ already from coverage.
In the zero branch both sides are $1$; in the branch $u=S(Pu)$ both
sides are $Pu$.  Thus $K$ and $W$ are provably equal on the whole domain,
even if they are distinct descriptions.  This equality does not appeal to
standard-model agreement or to uniqueness of an arbitrary recurrence.

\subsection{Four elementary uses of equation induction}

Choose descriptions for addition and truncated subtraction, with recursion
on the second argument:
\begin{align}
 x+0&=x, &x+Sy&=S(x+y),\label{eq:eic-addition}\\
 d(x,0)&=x, &d(x,Sy)&=P(d(x,y)).\label{eq:eic-subtraction}
\end{align}
Using these equations and \eqref{eq:eic-induction}, we now derive
\begin{align}
 0+y&=y,\label{eq:eic-left-zero}\\
 (Sx)+y&=S(x+y),\label{eq:eic-left-successor}\\
 d(Sx,Sy)&=d(x,y),\label{eq:eic-cancel-successors}\\
 d(Sa,a)&=1.\label{eq:eic-diagonal}
\end{align}
Each induction below has a single equation as its induction formula.

For \eqref{eq:eic-left-zero}, induct on $y$.  The base is $0+0=0$.
The step is $0+Sy=S(0+y)=Sy$.

For \eqref{eq:eic-left-successor}, hold $x$ fixed and induct on $y$.
At zero both sides are $Sx$.  At a successor the induction hypothesis
gives
\[
 (Sx)+Sy=S\bigl((Sx)+y\bigr)
       =S\bigl(S(x+y)\bigr)=S(x+Sy).
\]
Combining \eqref{eq:eic-left-successor} with the defining equation for
addition gives the useful rearrangement
\begin{equation}
 (Sy)+r=y+Sr.
 \label{eq:eic-shift-successor}
\end{equation}
Indeed, both sides are $S(y+r)$.

For \eqref{eq:eic-cancel-successors}, hold $x$ fixed and induct on $y$.
The base computes as
\[
 d(Sx,S0)=P(d(Sx,0))=P(Sx)=x=d(x,0).
\]
Applying $P$ to the induction equation gives the step:
\[
 d(Sx,S(Sy))=P(d(Sx,Sy))=P(d(x,y))=d(x,Sy).
\]

Finally, induct on $a$ for \eqref{eq:eic-diagonal}.  The base is
$d(S0,0)=S0=1$.  At a successor,
\[
 d(S(Sa),Sa)=d(Sa,a)=1,
\]
using \eqref{eq:eic-cancel-successors} and the induction hypothesis.

\subsection{An equation invariant that reflects actual equality}

The remaining arithmetic ingredient is the following equation, valid for
all $x,y$:
\begin{equation}
 C\bigl(\operatorname{sg}(d(x,y)),\ x,\ y+d(x,y)\bigr)=x.
 \label{eq:eic-residual-invariant}
\end{equation}
Its meaning is that a nonzero subtraction residual reconstructs the
original argument by addition.  The selector expresses this assertion as
one literal equation, so equation induction applies directly.

Hold $x$ fixed and induct on $y$.  At $y=0$, subtraction gives residual
$x$, and \eqref{eq:eic-left-zero} gives $0+x=x$.  The desired equation
is consequently $C(\operatorname{sg}(x),x,x)=x$, which is
\eqref{eq:eic-selector-idempotence}.

For the step, assume \eqref{eq:eic-residual-invariant} at $y$ and
abbreviate
\[
 r=d(x,y),\qquad r'=Pr=d(x,Sy).
\]
These are local term abbreviations, not additional function symbols.
Split coverage for $r'$.

If $r'=0$, then $\operatorname{sg}(r')=0$, and the successor instance
of \eqref{eq:eic-residual-invariant} reduces immediately to $x=x$.

In the other branch, $r'=S(Pr')$, so constructor separation gives
$r'\ne0$.  Split coverage for $r$.  Its zero branch would imply
$r'=P0=0$, a contradiction, from which the desired step equation follows.
In its successor branch,
\[
 r=S(Pr)=Sr',\qquad \operatorname{sg}(r)=1.
\]
The induction equation therefore computes to $y+r=x$.  Meanwhile,
\eqref{eq:eic-shift-successor} gives
\[
 (Sy)+r'=y+Sr'=y+r=x.
\]
Since $r'$ is a successor, $\operatorname{sg}(r')=1$, so this is
exactly the branch selected in the successor instance of
\eqref{eq:eic-residual-invariant}.  This completes the equation induction.
All case distinctions came from coverage; no induction on an implication
or a disjunction was used.

\paragraph{Reflection at residual $1$.}
We have, for arbitrary $a,b$,
\begin{equation}
 d(Sa,b)=1\quad\longleftrightarrow\quad a=b.
 \label{eq:eic-residual-reflection}
\end{equation}
For the forward direction, substitute $x=Sa$, $y=b$ into
\eqref{eq:eic-residual-invariant}.  If $d(Sa,b)=1$, the sign is $1$
and the selector returns its second branch, giving
\[
 b+1=Sa.
\]
The defining equation for addition gives $b+1=S(b+0)=Sb$.  Thus
$Sb=Sa$, and applying $P$ gives $b=a$, hence $a=b$.
Conversely, $a=b$ and \eqref{eq:eic-diagonal} give
$d(Sa,b)=d(Sa,a)=1$.

Define the equality characteristic by the fixed description
\begin{equation}
 e(a,b)=\operatorname{sg}\bigl(K(d(Sa,b))\bigr).
 \label{eq:eic-equality-characteristic}
\end{equation}
The sign, countdown, and residual reflection lemmas give
\begin{align}
 e(a,b)=0&\quad\longleftrightarrow\quad a=b,
       \label{eq:eic-equality-reflection}\\
 e(a,b)=0&\ \lor\ e(a,b)=1.
       \label{eq:eic-equality-range}
\end{align}
In detail, $e(a,b)=0$ is equivalent to $K(d(Sa,b))=0$ by
\eqref{eq:eic-sign-zero}, which is equivalent to $d(Sa,b)=1$ by
\eqref{eq:eic-countdown-zero}, and then to $a=b$ by
\eqref{eq:eic-residual-reflection}.  The range assertion is
\eqref{eq:eic-sign-range} applied to $K(d(Sa,b))$.

The reflection property concerns the logic's actual equality throughout
the domain.  It was not assumed as part of a numerical comparison symbol,
and its proof did not restrict $a$ or $b$ to standard numerals.
It also yields decidable equality.  In the $e(a,b)=0$ branch,
\eqref{eq:eic-equality-reflection} gives $a=b$.  In the $e(a,b)=1$
branch, an assumption $a=b$ would give $e(a,b)=0$ and therefore $0=1$,
contradicting separation.  Hence
\begin{equation}
 \forall a\,\forall b\,\bigl(a=b\lor(a=b\to\bot)\bigr)
 \label{eq:eic-equality-decidable}
\end{equation}
is a theorem, not an extra hypothesis.

\subsection{Compilation of all quantifier-free formulas}

For each fixed quantifier-free formula $A(\vec p)$, construct a term
$b_A(\vec p)$ by external recursion on its finite syntax.  Take equality,
$\bot$, $\land$, $\lor$, and $\to$ as the primitive formula constructors;
negation is implication into $\bot$.  Set
\begin{align}
 b_{s=t}&=e(s,t), & b_{\bot}&=1,\label{eq:eic-boolean-atoms}\\
 b_{A\land D}&=C(b_A,b_D,1),\label{eq:eic-boolean-and}\\
 b_{A\lor D}&=C(b_A,0,b_D),\label{eq:eic-boolean-or}\\
 b_{A\to D}&=C(b_A,b_D,0).\label{eq:eic-boolean-implies}
\end{align}
Every term here is a finite composition of fixed primitive-recursive
descriptions; if desired, projections and composition give it a single
description symbol.  No internal quantification over formulas or programs
is involved.

An external induction on the syntax of $A$ proves, within the arithmetic
theory,
\begin{equation}
 \bigl(b_A=0\lor b_A=1\bigr)
 \quad\text{and}\quad
 \bigl(b_A=0\longleftrightarrow A\bigr).
 \label{eq:eic-boolean-correctness}
\end{equation}
Here are the full connective checks.  In every case $b_A=1$ implies
$\neg A$: if $A$ held, the already established implication
$A\to b_A=0$ would contradict separation.

\paragraph{Atomic formulas and falsity.}
For an equation, both claims are
\eqref{eq:eic-equality-reflection}--\eqref{eq:eic-equality-range}.
For falsity the range assertion is $1=0\lor1=1$, with its right disjunct
proved by reflexivity.  Also $1=0\to\bot$ follows from separation and
symmetry, while $\bot\to1=0$ follows by explosion.  Thus $1=0$ is
equivalent to $\bot$, as required.

\paragraph{Conjunction.}
Split $b_A=0\lor b_A=1$.  In the zero branch, $A$ holds and
$b_{A\land D}=b_D$.  Its range follows from the range of $b_D$,
and its being zero is equivalent to $D$, hence to $A\land D$ under
the known $A$.  In the one branch, $\neg A$ holds and
$b_{A\land D}=1$.  Its range is immediate.  If it were zero,
separation would give $\bot$, hence $A\land D$; conversely an
assumption $A\land D$ contradicts $\neg A$, and explosion gives
$b_{A\land D}=0$.

\paragraph{Disjunction.}
In the branch $b_A=0$, the formula $A$ holds and
$b_{A\lor D}=0$.  The range is immediate, and both $b_{A\lor D}=0$
and $A\lor D$ hold, proving their equivalence.  In the branch
$b_A=1$, we have $\neg A$ and $b_{A\lor D}=b_D$.  Its range follows
from the range of $b_D$.  If its value is zero, then $D$ and hence
$A\lor D$.  Conversely, from $A\lor D$, the $A$ alternative
contradicts $\neg A$ and yields $b_D=0$ by explosion, while the
$D$ alternative yields $b_D=0$ directly.  This proves the required
equivalence in both alternatives.

\paragraph{Implication.}
In the branch $b_A=0$, the formula $A$ holds and
$b_{A\to D}=b_D$.  Its range follows from the range of $b_D$.
If its value is zero, then $D$, which gives $A\to D$.  Conversely,
$A\to D$ and the known $A$ give $D$, hence $b_D=0$.
In the branch $b_A=1$, we have $\neg A$ and
$b_{A\to D}=0$.  Its range is immediate.  To prove $A\to D$,
assume $A$; this contradicts $\neg A$, and explosion gives $D$.
Thus both $b_{A\to D}=0$ and $A\to D$ hold in this branch, and
their equivalence follows.  This argument uses the proved two-value
range of the finite characteristic.  It does not use excluded middle
for arbitrary formulas or change the surrounding intuitionistic logic.

This completes the external syntactic induction proving
\eqref{eq:eic-boolean-correctness} for every quantifier-free formula.

\subsection{Recovery of induction and the full theory}

Let $A(n,\vec p)$ be any fixed quantifier-free formula.  Suppose an
arbitrary retained context $\Gamma$ proves its base and its induction
step:
\[
 \Gamma\vdash A(0,\vec p),
 \qquad
 \Gamma,A(z,\vec p)\vdash A(Sz,\vec p),
\]
with the usual freshness condition on $z$.  By
\eqref{eq:eic-boolean-correctness}, the base gives
$\Gamma\vdash b_A(0,\vec p)=0$.  Under the additional assumption
$b_A(z,\vec p)=0$, the same equivalence gives $A(z,\vec p)$,
the given step gives $A(Sz,\vec p)$, and the equivalence gives
$b_A(Sz,\vec p)=0$.  Apply literal-equation induction to
$b_A(n,\vec p)=0$ and then apply its reflection implication.
The result is $\Gamma\vdash A(t,\vec p)$ for every target term $t$.
Thus the theory proves the full contextual quantifier-free induction
rule, including in arbitrary quantified contexts.

\paragraph{Completion theorem.}
In the common language of fixed primitive-recursive descriptions, the
theory
\[
 \mathsf{PRE}+\mathsf{Cov}+\mathsf{Sep}+\mathsf{EI}
\]
has exactly the same sequents as ordinary first-order intuitionistic
primitive recursive arithmetic with all quantifier-free induction
instances.  The forward simulation follows from the recovered
quantifier-free induction, successor injectivity, and successor separation;
all primitive-recursive defining equations and all logical rules are
shared.  For the reverse simulation, literal-equation induction is an
instance of quantifier-free induction, and successor separation gives
$\mathsf{Sep}$.  Quantifier-free induction proves $\mathsf{Cov}$:
at zero choose its left disjunct, and at a successor use
$Sx=S(P(Sx))$, which follows from $P(Sx)=x$, for its right disjunct.

Both interpretations therefore use the full domain, defined by $x=x$,
and preserve every function symbol, equality, logical connective, and
quantifier.  Both composites have the identity isomorphism, with graph
$x=y$.  The theories are consequently identity-domain bi-interpretable
and equiconsistent.  This completion theorem is unconditional once the
literal-equation induction rule has been obtained; it has no dependency
on a classical successor-uniqueness reduction.

\section{Pairing, primitive recursion, and internal numerical composites}
\label{sec:iterator-coding}
This appendix is adapted with attribution from the progress-closure
arithmetic development~\cite{progress}. Its prerequisites have now been
proved: scalar predecessor, arithmetic reflection, and full quantifier-free
induction. Pairing is a derived proof tool; it is not part of any checker.

\subsection{Pairing from a bounded square-root search}

Once QFI is available, a familiar square-based injection gives a
particularly direct compiler. Write $a\le b$ for $d(a,b)=0$.
The usual linear-order and semiring identities follow by fixed
quantifier-free inductions from the successor, addition, and monus
equations. Multiplication itself is the scalar term
\[
                        a\cdot b=I(b,0;u.u+a).
\]
Only primitive-recursive arithmetic is involved in the following
bounded root search; there is no minimization over an unbounded domain.

For a fixed parameter $n$, put
\[
 U_n(v)=C(d((Sv)^2,n),Sv,v),\qquad
 R_n(i)=I(i,0;v.U_n(v)),\qquad q(n)=R_n(n).
\]
The update increments $v$ exactly when $(v+1)^2\le n$, and
otherwise retains it. QFI in $i$, with $n$ fixed, proves
\[
 R_n(i)\le i,\qquad R_n(i)^2\le n,\qquad
 R_n(i)=i\ \vee\ n<(S(R_n(i)))^2.              \tag{Root invariant}
\]
The base has $R_n(0)=0$. If the update increments, its test
contradicts the strict-upper-bound alternative in the third invariant.
Thus the old state equals $i$, and the new state equals $Si$;
its square bound is the increment test. If the update does not increment, the
strict upper square bound holds and the state remains fixed.
An already established strict upper square bound forces every
subsequent update to remain fixed. These are exactly the cases
needed for the three displayed invariants.

At $i=n$, either the strict upper square bound already holds or
$R_n(n)=n$, in which case $n<(n+1)^2$ supplies it. Thus
\[
                        q(n)^2\le n<(q(n)+1)^2.           \tag{Root}
\]
Monotonicity of squares shows uniqueness of a number satisfying
Root. All formulas just used are quantifier free in the scalar
language, with fixed $n$ during the induction.

Now define
\[
 \pi(a,b)=(a+b)^2+a,\qquad
 L(n)=d(n,q(n)^2),\qquad H(n)=d(q(n),L(n)).
\]
For $k=a+b$, one has $k^2\le k^2+a<(k+1)^2$ since $a\le k$.
Root uniqueness gives $q(\pi(a,b))=a+b$, and therefore
\[
                    L(\pi(a,b))=a,\qquad H(\pi(a,b))=b.
\]
Surjectivity of $\pi$ is unnecessary. The original numerical domain
is never replaced by the range of this coding; pairing is used only
to store the finite state of a fixed primitive recursion. This
standard square pairing is a convenient reconstruction here, not
a claimed new arithmetic basis.

\subsection{Compiling each primitive recursion}

Suppose the fixed target description is
\[
 f(0,\vec p)=g(\vec p),\qquad
 f(Sn,\vec p)=h(n,f(n,\vec p),\vec p).
\]
Assume by external structural induction that the earlier $g,h$
already have scalar translations $g^*,h^*$. Define
\[
 W(u,\vec p)=\pi(SL(u),h^*(L(u),H(u),\vec p)),
\]
\[
 V(n,\vec p)=I(n,\pi(0,g^*(\vec p));u.W(u,\vec p)),
 \qquad f^*(n,\vec p)=H(V(n,\vec p)).
\]
The update has only its state variable bound, with every $p_i$
fixed. Equation induction proves $L(V(n,\vec p))=n$.
The two projection equations and scalar computation then prove
both defining equations for $f^*$ internally. Initial functions
and composition translate directly. This completes the translation
of every fixed finite PR description without adding a universal
evaluator.

\subsection{Why both numerical round trips are internal}

Two fixed scalar iterators with the same seed and provably equal
updates are internally equal at all counters: use literal-equation
induction comparing their values. Their base agrees, and their
successor equations plus congruence prove the step. The same
uniqueness argument works for ordinary primitive recursion.

For an ordinary PR description, induct externally on its finite
construction. The double translation satisfies its original
defining equations, so internal equation induction identifies it
with the original operation. For a native scalar description,
earlier translated updates agree by the external induction
hypothesis; scalar uniqueness identifies the double-translated
iterator. These proofs hold at arbitrary internal counters.

Thus the domain relation $x=y$ is provably a total bijection
preserving every fixed numerical operation in both composites.
The formula constructions and round trips in
Section~\ref{sec:choice-introduction} give the corresponding formula
composites. Both interpretations preserve falsity, taking native
$0=1$ to an ordinary contradiction and ordinary $\bot$ to native
$0=1$. Finite proof translations therefore establish equiconsistency.


\clearpage
\begin{thebibliography}{9}
\raggedright
\bibitem{readme} V. Reshetnikov, ProveIt, ConciseIPRA README, pinned commit
\texttt{e1d9e3b12edd642ea2f5b572bb43febceaa50d93}.
\href{https://github.com/VladimirReshetnikov/ProveIt/blob/e1d9e3b12edd642ea2f5b572bb43febceaa50d93/Logic/PrimitiveRecursiveArithmetic/ConciseIPRA/README.md}{Repository README}.
\bibitem{six} ProveIt, \texttt{Lean/ConciseIPRA/NativeSix.lean} and
\texttt{NativeDiagonal.lean}, same pinned commit.
\href{https://github.com/VladimirReshetnikov/ProveIt/blob/e1d9e3b12edd642ea2f5b572bb43febceaa50d93/Logic/PrimitiveRecursiveArithmetic/ConciseIPRA/Lean/ConciseIPRA/NativeSix.lean}{Six-rule calculus and its theorems}.
\bibitem{mirror} ProveIt, \texttt{Research/mirror\_guard}, same commit.
\href{https://github.com/VladimirReshetnikov/ProveIt/tree/e1d9e3b12edd642ea2f5b572bb43febceaa50d93/Logic/PrimitiveRecursiveArithmetic/ConciseIPRA/Research/mirror_guard}{One-sided computation guard package}.
\bibitem{reflexive} ProveIt, \texttt{Research/reflexive\_base/proof.tex},
same commit.
\href{https://github.com/VladimirReshetnikov/ProveIt/blob/e1d9e3b12edd642ea2f5b572bb43febceaa50d93/Logic/PrimitiveRecursiveArithmetic/ConciseIPRA/Research/reflexive_base/proof.tex}{Reflexive-base introduction and substitution repair}.
\bibitem{bridge} ProveIt, \texttt{Lean/ConciseIPRA/NativeBiInterpretation.lean},
same commit.
\href{https://github.com/VladimirReshetnikov/ProveIt/blob/e1d9e3b12edd642ea2f5b572bb43febceaa50d93/Logic/PrimitiveRecursiveArithmetic/ConciseIPRA/Lean/ConciseIPRA/NativeBiInterpretation.lean}{Faithful readings, both round trips, and equiconsistency}.
\bibitem{erasure} ProveIt, \texttt{Research/certified\_zero/reference/}\allowbreak
\texttt{scalar\_predecessor.tex}, same commit. Reproduced in the arithmetic appendix.
\bibitem{forward} ProveIt, \texttt{Research/certified\_zero/reference/}\allowbreak
\texttt{equality\_reflection.tex}, same commit. Reproduced in the arithmetic appendix.
\bibitem{progress} ProveIt, \texttt{Research/certified\_zero/reference/}\allowbreak
\texttt{primitive\_recursion.tex}, same commit. Reproduced in the arithmetic appendix.
\end{thebibliography}
\end{document}
